\documentclass[../main.tex]{subfiles}
\begin{document}

\chapter{Virtual Memory}
\lessoninfo{
  video={Lesson 13, videos 1--27},
  textbook={H\&P \cite{hennessy2017} §B.4--B.5},
  keywords={virtual memory, physical memory, page, page frame, page table, page fault, virtual page number, page offset, multi-level page table, page size, internal fragmentation, TLB, TLB miss, page table walk},
}

Lesson 12 treated every address produced by the processor as a location in
main memory.  In a real system the addresses that programs use are
\emph{virtual}: every program sees a large address space of its own, and the
operating system and the processor together map the parts of it that the
program actually uses onto the memory that is installed.  This lesson explains
why virtual memory is needed, how page tables translate virtual addresses into
physical ones, how large these tables are and how multi-level page tables keep
them small, and why translation would slow down every memory access if the
processor did not cache recent translations in a translation lookaside buffer,
which Lesson 14 combines with the caches of Lesson 12.

\section{Why virtual memory}
\label{sec:l13-why}

Hardware and software see memory differently.\source{Lesson 13, videos 1--3;
H\&P §B.4; \cite{kilburn1962}.}  To the hardware, memory is what is installed
in the machine: an array of as many bytes as the memory modules provide.  To a
program, memory is an address space of $2^{32}$ bytes on a 32-bit processor or
$2^{64}$ bytes on a 64-bit processor, and all of it belongs to that program:
code, data and stack are placed at whatever addresses the program and its
compiler choose, regardless of how much memory is installed and of which other
programs are running at the same time.

The two views disagree in two respects.
\begin{itemize}
  \item The address space of a single program is much larger than the
    installed memory.  A 32-bit address space is 4~GB, which a machine may or
    may not have; a 64-bit address space is 16~EB (1~EB $= 2^{60}$ bytes),
    many orders of magnitude more than any machine has.\caution{No processor translates all 64 bits.  x86-64 requires
    \emph{canonical} addresses, in which bits 63--48 repeat bit~47, so the
    usable space is 256~TB rather than 16~EB.}
  \item Many programs run at the same time, and each of them uses the same
    addresses, as if it owned the entire address space.
\end{itemize}

\begin{definition}[Virtual memory]
  The mechanism that gives every program the view of a large memory of its
  own, and maps the parts of that memory the program actually uses onto the
  memory installed in the machine, is called \term{virtual memory}[仮想記憶].
\end{definition}

The operating system decides where each part of each program is kept, and the
processor applies these decisions to every address on every memory access.
Programs can thus be written without knowing the size of memory or the other
programs they will run with, a program cannot access the memory of another
unless the operating system arranges it, and together the programs can use
more memory than is installed, the excess being kept on disk
(\cref{sec:l13-page-faults}).\sidenote{The idea is
that old: the Manchester Atlas of 1962 presented its 16K words of core store
and 96K words of magnetic drum as one address space of 512-word pages, and
translated addresses through 32 associative \emph{page address registers} ---
the ancestor of the TLB of \cref{sec:l13-tlb}.}

\section{Physical memory and virtual address spaces}
\label{sec:l13-views}

The two views correspond to two kinds of addresses.\source{Lesson 13, videos
4--5; H\&P §B.4.}

\begin{definition}[Physical memory]
  The memory actually installed in the machine is its
  \term{physical memory}[物理メモリ].  Its bytes are numbered consecutively,
  from 0 up to the size of physical memory minus one, and the number that
  identifies a byte is its \term{physical address}[物理アドレス].
\end{definition}

The size of physical memory is determined by the installed modules, not by
the width of the processor's addresses.  A 32-bit machine can address 4~GB
but may have only 2~GB installed; a 64-bit machine has perhaps tens of
gigabytes, a tiny fraction of the $2^{64}$ bytes it could address.  In the
examples of this lesson a physical address has just as many bits as the
installed memory requires: 31 bits for 2~GB.

\begin{definition}[Virtual address space]
  The \term{virtual address space}[仮想アドレス空間] of a program is the full
  range of addresses the program can use, from 0 to $2^{n}-1$ on an $n$-bit
  processor, independently of physical memory.  An address in it is a
  \term{virtual address}[仮想アドレス].
\end{definition}

All addresses a program works with are virtual: the program counter, the
pointers stored in its data, and the addresses the processor computes for
loads and stores.  A program uses only a few regions of its virtual address
space (\cref{fig:l13-mapping}, left and right):
\begin{itemize}
  \item its \emph{code} and its global \emph{data} at low addresses;
  \item above them the \emph{heap}, which grows toward higher addresses as the
    program allocates memory dynamically;
  \item the \emph{stack} at the top of the address space, which grows toward
    lower addresses.
\end{itemize}
Between the heap and the stack lies a huge range of addresses that the
program never uses.  Every running program has its own virtual address space
laid out in this way, so different programs use the same virtual addresses
for entirely different contents.

\section{Pages and page tables}
\label{sec:l13-pages}

A mapping recorded separately for every virtual address would need as many
table entries as there are addresses.\source{Lesson 13, videos 6--7; H\&P
§B.4.}  Virtual memory therefore maps fixed-size ranges of addresses, just as
a cache moves blocks rather than individual bytes.

\begin{definition}[Page, page frame]
  The virtual address space is divided into consecutive ranges of equal size,
  each called a \term{page}[ページ].  Physical memory is divided into ranges of
  the same size, each called a \term{page frame}[ページフレーム], or frame for
  short.  The size is a power of two, typically 4~KB, and any page fits
  exactly into any frame.
\end{definition}

\begin{definition}[Page table]
  The \term{page table}[ページテーブル] of a program records, for every page
  of its virtual address space, whether the page is in physical memory and, if
  it is, which frame holds it.
\end{definition}

The operating system maintains one page table for every running program and
updates it whenever it places one of the program's pages into a frame or
removes one.  The page table of the running program is found through a
register set by the operating system and is read to translate the program's
addresses, usually by the processor itself and in some systems by the
operating system (\cref{sec:l13-tlb}).\margin{That register is CR3 on x86-64 and
TTBR0/TTBR1 on AArch64; a context switch reloads it, and the whole mapping
changes at once.}  \Cref{fig:l13-mapping} illustrates
the resulting mapping for two programs.
\begin{itemize}
  \item Any page may be placed in any frame.  Pages that are adjacent in a
    virtual address space need not be adjacent in physical memory, and the
    frames of different programs are interleaved.
  \item Pages that a program never uses are not mapped and occupy no frame;
    only the used regions of the address space consume physical memory.
  \item The page tables of two programs may map pages to the same frame.  This
    lets programs share memory, for example a single copy of the code of a
    program that is running twice.
  \item Not every mapped page has to be in physical memory; the others are on
    disk (\cref{sec:l13-page-faults}).
\end{itemize}

\begin{figure}[htbp]
  \centering
  % The virtual address spaces of two programs mapped onto the frames of
% physical memory; one page of each program is on disk, and the code page
% is shared.  Page 0 (low addresses) is drawn at the bottom; page numbers
% are shown at the left inside each page.
\begin{tikzpicture}[x=1cm,y=1cm, font=\footnotesize\sffamily, >={Stealth[length=1.8mm]},
  pg/.style={draw, minimum width=1.5cm, minimum height=0.42cm, inner sep=0pt, font=\scriptsize\sffamily},
  num/.style={font=\scriptsize\sffamily, text=ln-muted, inner sep=1pt},
  ttl/.style={font=\scriptsize\sffamily\bfseries, text width=2.6cm, align=center, anchor=south, inner sep=1pt}]
  \def\h{0.42}
  \colorlet{lnpone}{ln-accent!15}
  \colorlet{lnptwo}{ln-tip!22}
  % ---------------- program 1 (left) and program 2 (right)
  \foreach \cx/\pcol in {0.75/lnpone, 9.05/lnptwo} {
    \node[pg, fill=\pcol] at (\cx,0.5*\h) {\iflangja{コード}{code}};
    \node[pg, fill=\pcol] at (\cx,1.5*\h) {\iflangja{データ}{data}};
    \node[pg, fill=\pcol] at (\cx,2.5*\h) {\iflangja{ヒープ}{heap}};
    \node[pg, fill=ln-box, minimum height=1.68cm, text=ln-muted] at (\cx,5*\h) {\iflangja{未使用}{unused}};
    \node[pg, fill=\pcol] at (\cx,7.5*\h) {\iflangja{スタック}{stack}};
    \foreach \i/\y in {0/0.5, 1/1.5, 2/2.5, {3--6}/6.5, 7/7.5} {
      \node[num, anchor=west] at (\cx-0.73,\y*\h) {\i};
    }
  }
  \node[ttl] at (0.75,8*\h+0.08) {\iflangja{プログラム 1 の\\仮想アドレス空間}{program 1:\\virtual address space}};
  \node[ttl] at (9.05,8*\h+0.08) {\iflangja{プログラム 2 の\\仮想アドレス空間}{program 2:\\virtual address space}};
  % ---------------- physical memory: frames 0..5 drawn from 2h to 8h
  \node[ttl] at (4.9,8*\h+0.08) {\iflangja{物理メモリ}{physical memory}};
  \node[pg, fill=lnpone] (f0) at (4.9,2.5*\h) {\iflangja{P1 データ}{P1 data}};
  \fill[lnpone] (4.15,3*\h) rectangle (4.9,4*\h);
  \fill[lnptwo] (4.9,3*\h) rectangle (5.65,4*\h);
  \node[pg] (f1) at (4.9,3.5*\h) {\iflangja{コード}{code}};
  \node[pg, fill=lnptwo] (f2) at (4.9,4.5*\h) {\iflangja{P2 データ}{P2 data}};
  \node[pg, text=ln-muted] (f3) at (4.9,5.5*\h) {\iflangja{空き}{free}};
  \node[pg, fill=lnpone] (f4) at (4.9,6.5*\h) {\iflangja{P1 スタック}{P1 stack}};
  \node[pg, fill=lnptwo] (f5) at (4.9,7.5*\h) {\iflangja{P2 ヒープ}{P2 heap}};
  % ---------------- disk below physical memory
  \draw[rounded corners=2pt, fill=ln-box] (3.45,-0.72) rectangle (6.35,0.12);
  \node[pg, fill=lnpone, minimum width=1.25cm] (d1) at (4.2,-0.3) {\iflangja{P1 ヒープ}{P1 heap}};
  \node[pg, fill=lnptwo, minimum width=1.25cm] (d2) at (5.6,-0.3) {\iflangja{P2 スタック}{P2 stack}};
  \node[num, anchor=north] at (4.9,-0.74) {\iflangja{ディスク}{disk}};
  % ---------------- mappings of program 1
  \draw[->] (1.5,0.5*\h) -- (f1.west);
  \draw[->] (1.5,1.5*\h) -- (f0.west);
  \draw[->] (1.5,7.5*\h) -- (f4.west);
  \draw[->, dashed] (0,2.5*\h) -- (-0.4,2.5*\h) -- (-0.4,-0.3) -- (d1.west);
  % ---------------- mappings of program 2
  \draw[->] (8.3,0.5*\h) -- (f1.east);
  \draw[->] (8.3,1.5*\h) -- (f2.east);
  \draw[->] (8.3,2.5*\h) -- (f5.east);
  \draw[->, dashed] (9.8,7.5*\h) -- (10.2,7.5*\h) -- (10.2,-0.3) -- (d2.east);
\end{tikzpicture}

  \caption{The virtual address spaces of two programs (pages 0--7, low
    addresses at the bottom) mapped onto six frames of physical memory.  Page~0
    of both programs is mapped to the same frame; each program has one page on
    disk (dashed); the unused pages 3--6 are not mapped.}
  \label{fig:l13-mapping}
\end{figure}

A page table describes \emph{every} page of the virtual address space, used
or not; in its simplest form (\cref{sec:l13-flat}) it has an entry for each.
The number of pages it describes is therefore set by the size of the virtual
address space and the size of a page; the size of physical memory only
determines how many frames an entry can refer to.\caution{A page table is
indexed by virtual page, not by frame.  With less physical memory there are
fewer frames, but the page table must still describe just as many pages.}

\section{Page faults}
\label{sec:l13-page-faults}

Together, the running programs usually map more pages than physical memory has
frames.\source{Lesson 13, video 8; H\&P §B.4.}  The pages that do not fit are
kept on disk, and physical memory holds the pages that have been used
recently.

\begin{definition}[Page fault]
  An access to a page that the page table does not map to a frame causes a
  \term{page fault}[ページフォールト]: an exception that transfers control to
  the operating system.
\end{definition}

If the page is on disk, the operating system handles the page fault in three
steps.
\begin{enumerate}
  \item It chooses a frame for the page.  If no frame is free, it evicts a page
    that has not been used recently; if that page has been written since it
    was loaded, it is first written back to disk.  The page table of the
    evicted page is updated to show that the page is no longer in memory.
  \item It reads the faulting page from disk into the frame and updates the
    page table to map the page to this frame.
  \item It restarts the instruction that caused the fault, which now finds its
    page in memory; this requires precise exceptions (Lesson 8).
\end{enumerate}

A disk access takes milliseconds, i.e.\ millions of clock cycles, so a page
fault is enormously more expensive than a cache miss.\margin{At \SI{3}{\giga\hertz} a
\SI{10}{\milli\second} disk seek is 30 million cycles; even an NVMe SSD read
of \SI{100}{\micro\second} is 300\,000.}  Page faults must be
rare, and it pays to let the operating system choose the page to evict with a
careful algorithm in software.

Not every page fault involves the disk: if the address lies outside every
region that the program has been allocated, there is nothing to load, and the
operating system treats the access as an error, usually by terminating the
program.\margin{A third case touches no disk: Linux counts a fault that
only has to allocate a zeroed frame or copy a shared page as a \emph{minor}
fault, against \emph{major} faults that read from storage.}

\begin{note}
  Physical memory acts as a cache of the pages on disk.  Pages correspond to
  the blocks of Lesson 12, frames to lines, page faults to cache misses, and
  writing only modified pages back to disk to a write-back cache with a dirty
  bit.  Placement is fully associative, since any page may be placed in any
  frame, yet no tags have to be searched, because the page table is indexed
  directly by the page.
\end{note}

\section{Address translation}
\label{sec:l13-translation}

Because pages start at multiples of the page size, a virtual address splits
into two fields, just like the block number and block offset of a cache
address.\source{Lesson 13, videos 9--10; H\&P §B.4.}

\begin{definition}[Virtual page number, page offset]
  With pages of $2^{k}$ bytes, the low-order $k$ bits of a virtual address form
  its \term{page offset}[ページオフセット], the position of the byte within
  its page.  The remaining high-order bits form its
  \term{virtual page number}[仮想ページ番号] (VPN)%
  \index{VPN|see{virtual page number}}, which identifies the page.
\end{definition}

Frames are numbered in the same way: the frame that starts at physical address
$j \times 2^{k}$ has the \term{physical frame number}[物理フレーム番号]
(PFN)\index{PFN|see{physical frame number}}~$j$.  The page table maps the VPN
to the PFN, and translation replaces the one by the other while keeping the
offset (\cref{fig:l13-translation}):
\begin{equation}
  \begin{aligned}
    \text{VPN} &= \bigl\lfloor \mathit{VA} / 2^{k} \bigr\rfloor , &
    \text{offset} &= \mathit{VA} \bmod 2^{k} , \\
    \text{PFN} &= \text{page table}[\text{VPN}] , &
    \mathit{PA} &= \text{PFN} \times 2^{k} + \text{offset} ,
  \end{aligned}
  \label{eq:l13-translation}
\end{equation}
where $\mathit{VA}$ and $\mathit{PA}$ are the virtual and physical addresses.
The offset passes through unchanged because a page and its frame have the same
size and the byte keeps its position within them; in binary, the physical
address is the PFN bits followed by the offset bits.\caution{Translation never
moves a byte out of its page, but an unaligned access can straddle a page
boundary: its two halves lie in different pages and need two translations,
into frames that may be far apart.}  The VPN has $k$ bits
fewer than a virtual address and the PFN $k$ bits fewer than a physical
address, so the two numbers differ in width whenever virtual and physical
addresses do.  If the page table holds no PFN for the page, the access causes a
page fault instead.

\begin{figure}[htbp]
  \centering
  % Translation with a flat page table: 32-bit virtual address, 4 KB pages,
% 30-bit physical address.  The VPN indexes the page table, the PFN found
% there replaces it, and the page offset is copied unchanged.
\begin{tikzpicture}[x=1cm,y=1cm, font=\footnotesize\sffamily, >={Stealth[length=1.8mm]},
  bits/.style={font=\scriptsize\sffamily, text=ln-muted, inner sep=1pt},
  cellt/.style={font=\scriptsize\ttfamily}]
  % --- virtual address
  \draw[fill=ln-accent!15] (0,0) rectangle (3.2,0.55);
  \node at (1.6,0.275) {VPN\enspace\texttt{0x0040A}};
  \draw[fill=white] (3.2,0) rectangle (5.6,0.55);
  \node at (4.4,0.275) {\iflangja{オフセット}{offset}\enspace\texttt{0x7C6}};
  \node[bits, anchor=south] at (1.6,0.55) {31\,\dots\,12};
  \node[bits, anchor=south] at (4.4,0.55) {11\,\dots\,0};
  \node[anchor=west] at (5.7,0.275) {\iflangja{仮想アドレス}{virtual address}};
  % --- page table: V | PFN, 6 drawn rows
  \node[font=\sffamily\scriptsize\bfseries, anchor=south] at (2.1,-0.92) {\iflangja{ページテーブル}{page table}};
  \node[bits, anchor=south] at (1.15,-1.2) {V};
  \node[bits, anchor=south] at (2.35,-1.2) {PFN};
  \fill[ln-accent!15] (0.9,-2.46) rectangle (3.3,-2.88);
  \foreach \r/\v/\pfn/\lab in {0/1/0x00A51/0, 1/0/{{\sffamily ---}}/1, 2/{}/{}/{}, 3/1/0x1F3C5/0x0040A, 4/{}/{}/{}, 5/1/0x3B207/0xFFFFF} {
    \draw (0.9,-1.2-0.42*\r) rectangle (1.4,-1.62-0.42*\r);
    \draw (1.4,-1.2-0.42*\r) rectangle (3.3,-1.62-0.42*\r);
    \node[cellt] at (1.15,-1.41-0.42*\r) {\v};
    \node[cellt] at (2.35,-1.41-0.42*\r) {\pfn};
    \node[bits, font=\scriptsize\ttfamily, anchor=west] at (3.35,-1.41-0.42*\r) {\lab};
  }
  \foreach \r in {2,4} {
    \node[font=\small, text=ln-muted, inner sep=0pt] at (1.15,-1.41-0.42*\r) {$\vdots$};
    \node[font=\small, text=ln-muted, inner sep=0pt] at (2.35,-1.41-0.42*\r) {$\vdots$};
  }
  % --- VPN selects an entry
  \draw[->] (0.3,0) -- (0.3,-2.67) -- (0.9,-2.67);
  % --- physical address
  \draw[->] (2.35,-3.72) -- (2.35,-4.3);
  \draw[fill=ln-accent!15] (0.8,-4.85) rectangle (3.2,-4.3);
  \node at (2.0,-4.575) {PFN\enspace\texttt{0x1F3C5}};
  \draw[fill=white] (3.2,-4.85) rectangle (5.6,-4.3);
  \node at (4.4,-4.575) {\iflangja{オフセット}{offset}\enspace\texttt{0x7C6}};
  \node[bits, anchor=north] at (2.0,-4.85) {29\,\dots\,12};
  \node[bits, anchor=north] at (4.4,-4.85) {11\,\dots\,0};
  \node[anchor=west] at (5.7,-4.575) {\iflangja{物理アドレス}{physical address}};
  % --- offset is copied
  \draw[->] (5.0,0) -- (5.0,-4.3);
\end{tikzpicture}

  \caption{Translation of the virtual address \texttt{0x0040A7C6} with 4~KB
    pages and 30-bit physical addresses (\cref{ex:l13-translation}).  The VPN
    selects a page table entry; its PFN replaces the VPN, and the offset is
    copied.  An entry with V~$=0$ would cause a page fault.}
  \label{fig:l13-translation}
\end{figure}

\begin{example}
  \label{ex:l13-translation}
  A 32-bit processor with 4~KB pages ($k = 12$) has 1~GB of physical memory,
  hence 30-bit physical addresses.  The page offset is bits 11--0, the VPN
  bits 31--12 (20 bits), and a PFN has 18 bits.  The virtual address
  \texttt{0x0040A7C6} has VPN \texttt{0x0040A} and offset \texttt{0x7C6}.  The
  page table of \cref{fig:l13-translation} maps page \texttt{0x0040A} to frame
  \texttt{0x1F3C5}, so the physical address is
  $\texttt{0x1F3C5} \times 2^{12} + \texttt{0x7C6} = \texttt{0x1F3C57C6}$.
  The virtual address \texttt{0x000017FF} lies in page \texttt{0x00001}, which
  the same page table marks as not in memory: an access to it causes a page
  fault.
\end{example}

\tip{With 4~KB pages and addresses written in hexadecimal, the last three
digits are the offset; translation replaces the digits in front of them.}

\begin{keypoint}
  Translation replaces the virtual page number by a physical frame number and
  keeps the page offset.  The page table, maintained by the operating system,
  supplies the frame number; if it holds none for the page, the access causes
  a page fault.
\end{keypoint}

\section{The size of a flat page table}
\label{sec:l13-flat}

The simplest page table is a single array.\source{Lesson 13, videos 11--12;
H\&P §B.4.}

\begin{definition}[Flat page table]
  A \term{flat page table}[単一レベルページテーブル] is an array with one
  entry for every virtual page, indexed directly by the VPN; its element for a
  page is the \term{page table entry}[ページテーブルエントリ]%
  \index{PTE|see{page table entry}} (PTE) of that page.
\end{definition}

A PTE holds the PFN of its page and a few status bits: a valid bit that is set
only if the page is in physical memory, and bits that record how the page may
be accessed (read, write, execute).\margin{An x86-64 PTE is 8 bytes: the PFN plus a
present bit, read/write, user/supervisor, accessed and dirty bits, and a
no-execute bit at the top.}  A PTE is typically about as large as an
address, 4 bytes on a 32-bit processor and 8 bytes on a 64-bit processor, which
leaves room for the PFN and the status bits.  Since every virtual page has an entry,
\begin{equation}
  \text{size of a flat page table}
  \;=\; \frac{\text{size of the virtual address space}}{\text{page size}}
        \times \text{PTE size} .
  \label{eq:l13-flat-size}
\end{equation}
Two quantities are absent from \cref{eq:l13-flat-size}.  The size of physical
memory only affects how many bits of an entry the PFN occupies.  The amount of
memory the program actually uses does not matter at all, because unused pages
have entries too.  Every running program needs a page table of this size.\tip{Work in
exponents: with an $n$-bit address space, $2^{k}$-byte pages and $2^{s}$-byte
entries the table is $2^{\,n-k+s}$ bytes --- subtract the page bits, add the
entry bits.}

\begin{example}
  \label{ex:l13-flat-size}
  With 32-bit virtual addresses, 4~KB pages and 4-byte entries, a flat page
  table has $2^{32}/2^{12} = 2^{20}$ entries and occupies
  $2^{20} \times 4$~bytes $= 4$~MB for each running program, even for a
  program that uses only a few kilobytes.  With 64-bit virtual addresses,
  4~KB pages and 8-byte entries it would have $2^{52}$ entries and occupy
  $2^{55}$~bytes (32~PB), more than any memory or disk holds.  Even if only 48
  of the 64 address bits are used, the flat page table occupies
  $2^{36} \times 8 = 2^{39}$~bytes, i.e.\ 512~GB per program.
\end{example}

\section{Multi-level page tables}
\label{sec:l13-multilevel}

A flat page table is large because it has entries for the whole virtual
address space, although programs use only the bottom and the top of
it.\source{Lesson 13, videos 13--17; H\&P §B.4--B.5.}  A multi-level page
table avoids storing entries for the unused range in between.

\begin{definition}[Multi-level page table]
  A \term{multi-level page table}[多段ページテーブル] is a tree of small
  tables.  The VPN is split into one field per level.  The first field indexes
  the outer page table, whose entries point to tables of the next level; the
  entries of the tables of the last level, the inner page tables, hold PFNs.  An
  entry whose part of the address space is not used at all is marked invalid,
  and no tables exist below it.
\end{definition}

With two levels, a virtual address is split into an outer index $p_1$, an
inner index $p_2$ and the page offset (\cref{fig:l13-two-level}), and the
translation takes two steps:
\begin{enumerate}
  \item entry $p_1$ of the outer page table gives the location of an inner
    page table;
  \item entry $p_2$ of that inner page table gives the PFN, which is combined
    with the offset as in \cref{eq:l13-translation}.
\end{enumerate}
If either entry is invalid, no PFN is found and the access causes a page
fault.  The tables are usually sized so that each fits exactly into one
page.

\begin{figure}[htbp]
  \centering
  % Two-level translation: 32-bit virtual address split into a 10-bit outer
% index p1, a 10-bit inner index p2 and a 12-bit offset.  Outer entry p1
% points to an inner table, whose entry p2 holds the PFN.
\begin{tikzpicture}[x=1cm,y=1cm, font=\footnotesize\sffamily, >={Stealth[length=1.8mm]},
  bits/.style={font=\scriptsize\sffamily, text=ln-muted, inner sep=1pt},
  cellt/.style={font=\scriptsize\ttfamily},
  ttl/.style={font=\sffamily\scriptsize\bfseries, anchor=south, inner sep=1pt}]
  % --- virtual address
  \draw[fill=ln-accent!15] (0,0) rectangle (2.4,0.55);
  \node at (1.2,0.275) {$p_1 = \texttt{1}$};
  \draw[fill=ln-box] (2.4,0) rectangle (4.8,0.55);
  \node at (3.6,0.275) {$p_2 = \texttt{0x00A}$};
  \draw[fill=white] (4.8,0) rectangle (7.2,0.55);
  \node at (6.0,0.275) {\iflangja{オフセット}{offset}\enspace\texttt{0x7C6}};
  \node[bits, anchor=south] at (1.2,0.55) {31\,\dots\,22};
  \node[bits, anchor=south] at (3.6,0.55) {21\,\dots\,12};
  \node[bits, anchor=south] at (6.0,0.55) {11\,\dots\,0};
  \node[anchor=west] at (7.3,0.275) {\iflangja{仮想アドレス}{virtual address}};
  % --- outer table: V | inner table
  \node[ttl] at (1.7,-0.92) {\iflangja{外側テーブル}{outer table}};
  \node[bits, anchor=south] at (1.1,-1.2) {V};
  \node[bits, anchor=south] at (1.9,-1.2) {\iflangja{内側}{inner}};
  \fill[ln-accent!15] (0.9,-1.62) rectangle (2.5,-2.04);
  \foreach \r/\v/\p/\lab in {0/1/A/0, 1/1/B/{}, 2/1/C/2, 3/0/{{\sffamily ---}}/3, 4/{}/{}/{}, 5/1/D/1023} {
    \draw (0.9,-1.2-0.42*\r) rectangle (1.3,-1.62-0.42*\r);
    \draw (1.3,-1.2-0.42*\r) rectangle (2.5,-1.62-0.42*\r);
    \node[cellt] at (1.1,-1.41-0.42*\r) {\v};
    \node[cellt] at (1.8,-1.41-0.42*\r) {\p};
    \node[bits, anchor=east] at (0.85,-1.41-0.42*\r) {\lab};
  }
  \node[font=\small, text=ln-muted, inner sep=0pt] at (1.1,-1.41-0.42*4) {$\vdots$};
  \node[font=\small, text=ln-muted, inner sep=0pt] at (1.9,-1.41-0.42*4) {$\vdots$};
  \draw[->] (0.15,0) -- (0.15,-1.83) -- (0.9,-1.83);
  % index of outer entry 1, above the p1 arrow that enters its row
  \node[bits, anchor=south east] at (0.85,-1.78) {1};
  % --- inner table B: V | PFN, top edge at the height of outer entry 1
  \node[ttl] at (4.35,-1.55) {\iflangja{内側テーブル B}{inner table B}};
  \node[bits, anchor=south] at (3.65,-1.83) {V};
  \node[bits, anchor=south] at (4.6,-1.83) {PFN};
  \fill[ln-accent!15] (3.4,-3.09) rectangle (5.3,-3.51);
  \foreach \r/\v/\pfn/\lab in {0/1/0x2A017/0, 1/0/{{\sffamily ---}}/1, 2/{}/{}/{}, 3/1/0x1F3C5/10, 4/{}/{}/{}, 5/1/0x0C4E2/1023} {
    \draw (3.4,-1.83-0.42*\r) rectangle (3.9,-2.25-0.42*\r);
    \draw (3.9,-1.83-0.42*\r) rectangle (5.3,-2.25-0.42*\r);
    \node[cellt] at (3.65,-2.04-0.42*\r) {\v};
    \node[cellt] at (4.6,-2.04-0.42*\r) {\pfn};
    \node[bits, anchor=east] at (3.35,-2.04-0.42*\r) {\lab};
  }
  \foreach \r in {2,4} {
    \node[font=\small, text=ln-muted, inner sep=0pt] at (3.65,-2.04-0.42*\r) {$\vdots$};
    \node[font=\small, text=ln-muted, inner sep=0pt] at (4.6,-2.04-0.42*\r) {$\vdots$};
  }
  % --- pointer from outer entry 1 to the start of inner table B
  \fill (2.3,-1.83) circle (1pt);
  \draw[->] (2.3,-1.83) -- (3.4,-1.83);
  % --- p2 selects an entry of B
  \draw[->] (4.4,0) -- (4.4,-0.35) -- (5.75,-0.35) -- (5.75,-3.30) -- (5.3,-3.30);
  % --- physical address
  \draw[->] (4.6,-4.35) -- (4.6,-5.0);
  \draw[fill=ln-accent!15] (2.4,-5.55) rectangle (4.8,-5.0);
  \node at (3.6,-5.275) {PFN\enspace\texttt{0x1F3C5}};
  \draw[fill=white] (4.8,-5.55) rectangle (7.2,-5.0);
  \node at (6.0,-5.275) {\iflangja{オフセット}{offset}\enspace\texttt{0x7C6}};
  \node[bits, anchor=north] at (3.6,-5.55) {29\,\dots\,12};
  \node[bits, anchor=north] at (6.0,-5.55) {11\,\dots\,0};
  \node[anchor=west] at (7.3,-5.275) {\iflangja{物理アドレス}{physical address}};
  % --- offset is copied
  \draw[->] (6.8,0) -- (6.8,-5.0);
\end{tikzpicture}

  \caption{Two-level translation of the virtual address \texttt{0x0040A7C6}
    (\cref{ex:l13-two-level}).  Outer entry~1 points to inner table~B.  Entries
    A, C and D point to inner tables that are not drawn; invalid outer entries,
    such as entry~3, have no inner table at all.}
  \label{fig:l13-two-level}
\end{figure}

\begin{example}
  \label{ex:l13-two-level}
  The 20-bit VPN of \cref{ex:l13-translation} is split into a 10-bit outer
  index (bits 31--22) and a 10-bit inner index (bits 21--12).  With 4-byte
  entries every table has $2^{10}$ entries and occupies 4~KB, exactly one page.  The
  virtual address \texttt{0x0040A7C6} has VPN
  $\texttt{0x0040A} = 0000000001\,0000001010_2$, hence $p_1 = 1$ and
  $p_2 = \texttt{0x00A} = 10$.  Outer entry~1 points to inner table~B, and
  entry~10 of B holds PFN \texttt{0x1F3C5}.  The physical address is
  \texttt{0x1F3C57C6}, as with the flat page table: a multi-level page table
  stores the same mapping in a different form.
\end{example}

\subsection{Size of a multi-level page table}

The outer page table always exists, but an inner page table exists only for a
part of the address space in which at least one page is used.  With the
organization of \cref{ex:l13-two-level}, each inner table covers $2^{10}$
pages, i.e.\ 4~MB of the virtual address space.

\begin{example}
  \label{ex:l13-two-level-size}
  A program uses 10~MB at the bottom of its 32-bit address space for code,
  data and heap (addresses \texttt{0x00000000}--\texttt{0x009FFFFF}) and 1~MB
  at the top for its stack (\texttt{0xFFF00000}--\texttt{0xFFFFFFFF}).  With
  the two-level page table of \cref{ex:l13-two-level}, the bottom 10~MB lie in
  the ranges of outer entries 0, 1 and 2, of 4~MB each, and the stack in the
  range of outer entry 1023.  The page table consists of the outer table and
  four inner tables, $5 \times 4~\text{KB} = 20$~KB, instead of the 4~MB of a
  flat page table.
\end{example}

A two-level page table is larger than a flat one only if every outer entry
needs an inner table, i.e.\ if at least one page is used in each of the 1024
ranges of 4~MB.  It then occupies
$4~\text{KB} + 1024 \times 4~\text{KB} \approx 4.004$~MB, more than the flat
table by the size of the outer table.  Real programs use a few contiguous
regions of their address space and need only a few inner tables.\sidenote{A
different answer to the same problem is the \emph{inverted} page table: one
entry per frame rather than one per page, located by hashing the VPN.  Its
size follows physical memory and does not grow with the address space, but
collisions cost a chain walk and shared pages need extra machinery.  PowerPC
and UltraSPARC used it (H\&P §B.4).}

\subsection{More levels for 64-bit addresses}

Two levels do not suffice for 64-bit address spaces.  A 4~KB table of 8-byte
entries has $2^{9}$ entries, so each level of tables that fit into one page
can consume only 9 bits of the VPN.  With 48-bit virtual addresses the VPN has
36 bits.  A two-level page table with 9-bit inner tables would need a 27-bit
outer index and an outer table of $2^{27} \times 8$~bytes $= 1$~GB for every
program.  Splitting the 36 bits into four 9-bit fields instead yields a
four-level page table in which every table fits into one page; x86-64
processors use this organization.  Each further level makes room for 9 more
bits of virtual address.\margin{x86-64 calls the four levels PML4, PDPT, PD
and PT.  Processors from Ice Lake on offer a fifth (\texttt{CR4.LA57}) for
57-bit addresses, i.e.\ 128~PB.}

\begin{example}
  \label{ex:l13-four-level}
  In such a four-level page table, a table of the last level covers $2^{9}$
  pages, i.e.\ 2~MB of the address space; a table of the third level covers
  $2^{9}$ times as much, 1~GB; a table of the second level 512~GB; and the
  single table of the first level all $2^{48}$~bytes (256~TB).  A program that
  uses 8~MB at the bottom and 8~MB at the top of its address space needs, for
  each end, four tables of the last level and one of each of levels 3 and 2,
  and all of them hang below the single first-level table:
  $1 + 2 \times (4 + 1 + 1) = 13$ tables, or $13 \times 4~\text{KB} = 52$~KB,
  whereas its flat page table would occupy 512~GB (\cref{ex:l13-flat-size}).
\end{example}

\begin{keypoint}
  A page table must describe the entire virtual address space, but a
  multi-level page table needs storage only for the parts that are used.  The
  price is one table read per level on every translation.
\end{keypoint}

\section{Choosing the page size}
\label{sec:l13-page-size}

How large pages should be is a compromise between two
effects.\source{Lesson 13, video 18; H\&P §B.4.}
\begin{itemize}
  \item \textbf{Larger pages} mean fewer pages: fewer PTEs and a smaller page
    table, and, as \cref{sec:l13-tlb} shows, more memory covered by each
    translation that the processor caches.
  \item \textbf{Smaller pages} waste less memory in partly used pages.
\end{itemize}
The \term{page size}[ページサイズ], the number of bytes in a page and in a
frame, must balance the two.

\begin{definition}[Internal fragmentation]
  Memory is mapped a whole page at a time.  A region whose size is not a
  multiple of the page size therefore leaves part of its last page unused, and
  no other program can use that part.  The memory lost in this way is
  \term{internal fragmentation}[内部断片化].
\end{definition}

A region of arbitrary size leaves on average half a page unused, so internal
fragmentation grows in proportion to the page size.\caution{Paging removes
\emph{external} fragmentation, the waste left between allocations of different
sizes, because any page fits any frame.  What it leaves is the internal kind,
inside the last page.}

\begin{example}
  \label{ex:l13-page-size}
  A program on a 32-bit processor with 4-byte PTEs has three regions (code and
  data, heap, stack), each ending in a partly used page.
  \Cref{tab:l13-page-size} compares three page sizes.  Going from 4~KB to 4~MB
  pages shrinks the flat page table by a factor of $2^{10}$, from 4~MB to
  4~KB, but raises the expected internal fragmentation by the same factor,
  from 6~KB to 6~MB.
\end{example}

\begin{table}[htbp]
  \centering\small
  \caption{Page table size and expected internal fragmentation for the program
    of \cref{ex:l13-page-size} (half a page per region).}
  \label{tab:l13-page-size}
  \begin{tabular}{@{}rrrr@{}}
    \toprule
    Page size & Pages & Flat page table & Internal fragmentation \\
    \midrule
    4~KB  & $2^{20}$ & 4~MB   & $3 \times 2~\text{KB} = 6$~KB \\
    64~KB & $2^{16}$ & 256~KB & $3 \times 32~\text{KB} = 96$~KB \\
    4~MB  & $2^{10}$ & 4~KB   & $3 \times 2~\text{MB} = 6$~MB \\
    \bottomrule
  \end{tabular}
\end{table}

Pages are much larger than cache blocks (Lesson 12), because every page needs
its own PTE and every page brought in from disk costs a disk access.  The usual
page size is 4~KB.\margin{4~KB has been the x86 page size since the 386 of
1985.  AArch64 instead lets the operating system choose a granule of 4, 16 or
64~KB.}  Processors also support much larger pages, such as 2~MB
and 1~GB on x86-64, which the operating system can use for large regions that
a program fills completely, so that internal fragmentation hardly matters.

\section{Memory access time with translation}
\label{sec:l13-access-time}

Every memory access a program makes, whether it fetches an instruction, loads
or stores, uses a virtual address, which must be translated before physical
memory can be accessed.\source{Lesson 13, videos 19--21; H\&P §B.4.}  The page
table is itself kept in memory, so translation consists of memory reads.  With
a two-level page table, the processor performs a load such as
\texttt{LW R1,4(R2)} in four steps.
\begin{enumerate}
  \item It adds 4 to R2 to obtain the virtual address and takes its VPN.
  \item It computes the physical address of the outer PTE (the outer table
    address plus $p_1$ times the PTE size) and reads the PTE.
  \item It computes the address of the inner PTE in the same way from the
    table address in the outer PTE and $p_2$, and reads the PTE.
  \item It appends the offset to the PFN and accesses the data at this
    physical address through the cache.
\end{enumerate}
Each PTE read is an ordinary memory access that goes through the caches and may
hit or miss; if PTEs are not kept in the cache, each costs a full main-memory
access.\margin{Besides the ordinary caches, x86 processors hold
upper-level entries in dedicated \emph{paging-structure caches}, so a walk
usually needs only the last-level read.  Intel SDM vol.~3.}  With a page table of $L$ levels, every access of the program turns
into $L + 1$ memory accesses.  If PTE reads have the same average memory
access time (AMAT) as data accesses,
\begin{equation}
  \text{access time with translation} \;\approx\; (L + 1) \times \text{AMAT} .
  \label{eq:l13-walk-time}
\end{equation}

\begin{example}
  \label{ex:l13-access-time}
  A cache has a hit time of 2 cycles, a miss rate of \SI{5}{\percent} and a
  miss penalty of 40 cycles, so $\text{AMAT} = 2 + 0.05 \times 40 = 4$ cycles
  (\cref{eq:l12-amat}).  Without translation a memory access takes 4 cycles on
  average; with a flat page table it takes $2 \times 4 = 8$ cycles, and with a
  four-level page table $5 \times 4 = 20$ cycles.  Individual accesses vary
  widely: a load whose three upper-level PTEs hit ($3 \times 2$ cycles), whose
  last-level PTE misses ($2 + 40$ cycles) and whose data hits (2 cycles) takes
  50 cycles, 25 times as long as the cache hit it would be without
  translation.
\end{example}

The few upper-level entries are read by every translation and are usually
cached; the numerous last-level entries miss more often.  Either way, without
a faster way to translate, virtual memory would slow down every instruction
fetch, load and store by a large factor.

\section{The translation lookaside buffer}
\label{sec:l13-tlb}

A program that accesses a page is likely to access the same page again soon,
so the translations it needs at any time are few and are used over and
over.\source{Lesson 13, videos 22--24; H\&P §B.4.}  The processor therefore
caches them.

\begin{definition}[Translation lookaside buffer]
  A \term{translation lookaside buffer}[TLB]%
  \index{TLB|see{translation lookaside buffer}} (TLB) is a small, fast cache of
  recent translations.  Each entry holds a valid bit, a tag taken from the VPN
  (the whole VPN in a fully associative TLB), the PFN to which that page is
  mapped, and the access bits of its PTE.  If the VPN of an address matches a
  valid entry, the PFN is taken from the entry and no page table is read.
\end{definition}

A TLB is a cache in the sense of Lesson 12, but it can be much smaller than
the data cache, for three reasons.
\begin{itemize}
  \item It stores only translations, not data; an entry takes a few bytes.
  \item One entry serves a whole page rather than a single cache block, so a
    few entries cover the memory a program is actively using.\tip{The memory a
    TLB maps at one time is its \emph{reach}: entries $\times$ page size.
    64 entries of 4~KB reach 256~KB; the same 64 entries of 2~MB pages reach
    128~MB.}
  \item It stores the final result of the translation, not the entries of the
    individual levels.  A TLB hit takes a single lookup, however many levels
    the page table has.
\end{itemize}

\begin{example}
  \label{ex:l13-tlb-size}
  A 128~KB data cache with 64-byte blocks has 2048 lines.  With 4~KB pages,
  $128/4 = 32$ TLB entries cover as much memory as the cache holds, and 2048
  are needed if every cached block lies in a different page; spatial locality
  keeps the actual need far below this upper bound.
\end{example}

\subsection{TLB misses}

A \term{TLB miss}[TLB ミス] occurs when no valid TLB entry matches the VPN of
an address.  The translation must then be obtained from the page table.

\begin{definition}[Page table walk]
  A \term{page table walk}[ページテーブルウォーク] reads the page table
  entries of an address level by level, from the outer page table down to the
  entry that holds the PFN, as described in \cref{sec:l13-multilevel}.
\end{definition}

After a TLB miss, the walk finds the PFN, the translation is written into the
TLB, replacing an existing entry if the TLB is full, and the access proceeds
(\cref{fig:l13-tlb-flow}).\caution{A TLB miss is not a page fault: the walk
usually finds a valid PTE.}

The TLB must agree with the page tables: when the operating system changes a
mapping, e.g.\ by evicting a page, it invalidates the TLB entry of the page,
and when the processor switches to another program, the entries of the
previous program are invalidated or tagged with a program identifier.\sidenote{That tag is a PCID on x86 and an ASID on Arm; with it a context
switch need not flush the TLB.  Invalidating a single mapping is the harder
case: on a multicore processor every core that may hold the entry must drop
it, so the operating system sends an interprocessor interrupt and waits for
the acknowledgements.  This \emph{TLB shootdown} makes unmapping memory far
more expensive than mapping it.}

\begin{figure}[htbp]
  \centering
  % Translation of one memory access: TLB hit, TLB miss with a page table
% walk, and page fault handled by the operating system.
\begin{tikzpicture}[x=1cm,y=1cm, font=\footnotesize\sffamily, >={Stealth[length=1.8mm]},
  stp/.style={draw, fill=ln-box, align=center, inner sep=3pt, minimum height=0.8cm, text width=3.0cm},
  os/.style={stp, fill=white},
  dec/.style={draw, diamond, aspect=2.4, fill=ln-accent!15, align=center, inner sep=1pt},
  lab/.style={font=\sffamily\scriptsize, text=ln-muted}]
  % --- column 1: TLB lookup and the access itself
  \node[stp] (split) at (0,0)
    {\iflangja{仮想アドレスを VPN と\\オフセットに分ける}{split the virtual address\\into VPN and offset}};
  \node[stp] (look) at (0,-1.3)
    {\iflangja{VPN で TLB を検索}{look up the VPN\\in the TLB}};
  \node[dec] (hit) at (0,-2.65) {\iflangja{TLB ヒット?}{TLB hit?}};
  \node[stp] (acc) at (0,-5.45)
    {\iflangja{物理アドレス $=$ PFN と\\オフセット; データに\\アクセス}{physical address $=$\\PFN and offset;\\access the data}};
  % --- column 2: page table walk
  \node[stp] (walk) at (3.9,-2.65)
    {\iflangja{ページテーブルウォーク\\（ハードウェアまたは OS）}{page table walk\\(hardware or OS)}};
  \node[dec] (valid) at (3.9,-4.0) {\iflangja{PTE 有効?}{PTE valid?}};
  \node[stp] (fill) at (3.9,-5.45)
    {\iflangja{VPN と PFN を\\TLB に書き込む}{write VPN and PFN\\into the TLB}};
  % --- column 3: page fault
  \node[os] (fault) at (7.8,-4.0)
    {\iflangja{ページフォールト（OS）:\\ページをディスクから\\読み込み，PTE を更新}{page fault (OS):\\load the page from disk,\\update the page table}};
  % --- arrows
  \draw[->] (split) -- (look);
  \draw[->] (look) -- (hit);
  \draw[->] (hit) -- node[lab, right, pos=0.12] {\iflangja{はい}{yes}} (acc);
  \draw[->] (hit) -- node[lab, above] {\iflangja{いいえ}{no}} (walk);
  \draw[->] (walk) -- (valid);
  \draw[->] (valid) -- node[lab, right, pos=0.2] {\iflangja{はい}{yes}} (fill);
  \draw[->] (fill) -- (acc);
  \draw[->] (valid) -- node[lab, above] {\iflangja{いいえ}{no}} (fault);
  \draw[->] (fault.north) -- (7.8,-1.3) -- node[lab, above, pos=0.5] {\iflangja{命令を再実行}{restart the instruction}} (look.east);
\end{tikzpicture}

  \caption{Translation on a memory access.  A TLB hit supplies the PFN at
    once; a TLB miss requires a page table walk; an invalid PTE causes a page
    fault; once the page is in memory, the instruction is restarted.}
  \label{fig:l13-tlb-flow}
\end{figure}

The walk can be performed in two ways.
\begin{description}
  \item[In software] the TLB miss raises an exception, and a handler in the
    operating system walks the page table and writes the translation into the
    TLB.  The hardware is simple and the page tables may be organized in any
    way, but every TLB miss pays for an exception.  MIPS processors handle TLB
    misses this way.
  \item[In hardware] the processor contains a page table walker, logic that
    reads the PTEs and fills the TLB itself without involving the operating
    system.  TLB misses are handled much faster, but the page table must have
    exactly the format that the hardware expects.  High-performance processors,
    such as those of the x86 family, walk the page table in hardware.
\end{description}

\needspace{6\baselineskip}
\begin{keypoint}
  Every memory access needs a translation.  The TLB makes the common case, a
  TLB hit, a single fast lookup; a page table walk, in hardware or by the
  operating system, is needed only on a TLB miss, and the page-fault handler
  only when the walk finds no PFN.
\end{keypoint}

\section{TLB organization and performance}
\label{sec:l13-tlb-org}

The organization of a TLB follows from its role.\source{Lesson 13, videos
25--26; H\&P §B.4.}
\begin{description}
  \item[Associativity] Every TLB miss costs a walk of several memory reads, so
    TLBs are fully associative or highly set-associative; for so few entries
    the comparators are affordable, and translations do not compete for
    entries.
  \item[Size] The TLB should hit about as often as the cache, so it must cover
    the pages of the data the cache holds (\cref{ex:l13-tlb-size}); this takes
    about 64 to 512 entries.  Consulted on every access, it must also be about
    as fast as an L1 cache hit, which keeps it from growing larger.\margin{Skylake, per Intel's optimization manual: 64 first-level data TLB
    entries for 4~KB pages, 128 for instructions, and a 1536-entry 12-way
    second-level TLB behind both.}
  \item[Levels] If that does not cover enough memory, a larger and slower
    second-level TLB with hundreds to thousands of entries is placed behind the
    first, as caches are organized in levels (Lesson 14).  It is looked up only
    on a miss in the first, and a walk is needed only if both TLBs miss.
\end{description}

\begin{example}
  \label{ex:l13-tlb-performance}
  The processor of \cref{ex:l13-access-time} (AMAT of 4 cycles for every
  memory read, four-level page table) is given a TLB that is looked up in
  1 cycle and misses on \SI{1}{\percent} of the accesses.  Every access pays
  for the lookup and for the data access, and one in a hundred also for a walk
  of $4 \times 4 = 16$ cycles, so an access takes on average
  $1 + 0.01 \times 16 + 4 = 5.16$ cycles instead of 20.  Translation now costs
  1.16 cycles instead of 16, mostly for the lookup, which Lesson 14 overlaps
  with the cache lookup.\margin{This has the form of \cref{eq:l12-amat}:
  lookup time $+$ miss rate $\times$ penalty, with the walk as the penalty.
  A TLB is a cache and is evaluated as one.}
\end{example}

\begin{keypoint}[Lesson summary]
  Virtual memory maps the pages of each program's virtual address space onto
  frames of physical memory and keeps the rest on disk; an access to a page
  without a frame is a page fault.  Translation replaces the virtual page
  number by a frame number from the page table and keeps the offset
  (\cref{eq:l13-translation}).
  Multi-level page tables avoid the entries a flat table spends on unused
  address space (\cref{eq:l13-flat-size}); the page size trades table size
  against internal fragmentation.  Every level adds a memory read to each
  access, so a small, highly associative TLB caches translations, and the
  page table is walked only on a TLB miss.
\end{keypoint}

\end{document}
